% Generated draft: requires mathematical and source review.
\documentclass[11pt]{article}
\usepackage[margin=25mm]{geometry}
\usepackage{fontspec}
\setmainfont{Latin Modern Roman}
\usepackage{amsmath,amssymb,mathtools,unicode-math}
\setmathfont{Latin Modern Math}
\usepackage{enumitem,needspace,xurl,hyperref}
\setlength{\parindent}{0pt}
\setlength{\parskip}{5pt}
\setlength{\emergencystretch}{3em}
\newcommand{\N}{\mathbb{N}}
\newcommand{\Z}{\mathbb{Z}}
\newcommand{\Q}{\mathbb{Q}}
\newcommand{\R}{\mathbb{R}}
\newcommand{\C}{\mathbb{C}}
\newcommand{\F}{\mathbb{F}}
\newcommand{\A}{\mathbb{A}}
\newcommand{\PP}{\mathbb{P}}
\newcommand{\E}{\mathbb{E}}
\newcommand{\eps}{\varepsilon}
\newcommand{\dd}{\,\mathrm{d}}
\begin{document}
\section*{UnsolvedMath ID: 701403}
\textbf{Source problem code:} catalogue entry
\subsection{Definitions and mathematical statement}
The source attributes to P. Hall the following definition: a group \(G\) is homogeneous if, for every pair of finitely generated subgroups \(H_1,H_2\leq G\), every isomorphism \(H_1\to H_2\) extends to an automorphism of \(G\). It also attributes to Higman--Neumann--Neumann the assertion that every group embeds into a homogeneous group. The question asks for an analogous assertion for representations of groups.

The coefficient system and the morphisms of the intended representation category are not specified in the supplied record. The following is an explicit formulation under the usual interpretation using a fixed coefficient ring and equivariant pairs of maps; these choices require confirmation.

Fix a unital coefficient ring \(K\). A representation \((V,G)\) consists of a group \(G\) and a unital left \(K[G]\)-module \(V\). A morphism \((f,\varphi):(V,G)\to(W,L)\) consists of a \(K\)-linear map \(f:V\to W\) and a group homomorphism \(\varphi:G\to L\) satisfying
\[
f(g\cdot v)=\varphi(g)\cdot f(v)\qquad(g\in G,\ v\in V).
\]
It is an embedding if both maps are injective, and an isomorphism if both are bijective. A subrepresentation \((U,H)\) of \((V,G)\) consists of a subgroup \(H\leq G\) and a \(K[H]\)-submodule \(U\subseteq V\). As specified in the source, \((U,H)\) is finitely generated precisely when \(H\) is finitely generated as a group and \(U\) is finitely generated as a \(K[H]\)-module.

Under this interpretation, call \((W,L)\) homogeneous if, for any finitely generated subrepresentations \((U_1,H_1)\) and \((U_2,H_2)\), every isomorphism
\[
(f,\varphi):(U_1,H_1)\longrightarrow(U_2,H_2)
\]
extends to an automorphism \((F,\Phi)\) of \((W,L)\), meaning
\[
F|_{U_1}=f,\qquad \Phi|_{H_1}=\varphi.
\]
The proposed representation analogue is the assertion
\[
\begin{gathered}
\text{For every group }G\text{ and every left }K[G]\text{-module }V,\\
\text{there exist a homogeneous representation }(W,L)
\text{ and an embedding }(V,G)\hookrightarrow(W,L).
\end{gathered}
\]
No finite-generation assumption is imposed on the input representation. The source does not determine whether \(K\) must be a field, what other coefficient rings are allowed, or how the question is quantified over coefficients. Consequently, the displayed assertion is a conditional formalization of the intended categorical analogue, pending resolution of those qualifications.
\subsection{Short English statement}
Does every representation of a group embed into a representation in which every isomorphism between finitely generated subrepresentations extends to an automorphism? Here finite generation requires both a finitely generated group and a finitely generated module over its group algebra. The coefficient system and intended category require clarification. The supplied registry labels the problem open and records unverified AI solution claims.
\subsection{Sources}
\begin{itemize}
\item Local UnsolvedMath registry, supplied record with permanent ID \(701403\), exact source code \texttt{catalogue entry}, title \emph{Kourovka Notebook Problem 17.101}; proposed by B. I. Plotkin in 2010.
\item Supplied source attribution: \emph{The Kourovka Notebook (Unsolved Problems in Group Theory)}, 17th issue (2010), Problem 17.101; also the 21st issue edition (2026), identified in the record as \texttt{arXiv:1401.0300v48}. The supplied collection description names E. I. Khukhro and V. D. Mazurov as editors. The problem text attributes the group-homogeneity definition to P. Hall and the group embedding assertion to Higman--Neumann--Neumann; no separate bibliographic references for those attributions are supplied.
\item Supplied discussion: Achyuth Jayadevan, \emph{Preprint of 11 September 2026}, \url{https://zenodo.org/records/22722081}. The record explicitly describes this as an unverified AI solution claim asserting an affirmative answer.
\item Supplied discussion: M. Kinyon and J. Urban, \emph{Preprint of 17 September 2026}, \url{https://kourovkanotebookorg.wordpress.com/wp-content/uploads/2026/09/kourovka-mathematics.pdf}. The record explicitly describes this as an unverified AI solution claim asserting an affirmative answer.
\item Supplied status review, checked on \texttt{2026-10-16}: canonical status \texttt{open}, evidence classification \texttt{unrefereed\_claim}. Its notes report a \(\#\) marker for an unverified AI solution claim in the October 2026 21st issue. This assessment is supplied dated curation, not independently verified source consultation or mathematical verification.
\end{itemize}
\end{document}
